\chapter{System Overview and Hardware Architecture}

\section{Requirements}

From the objectives in Chapter~1, the hardware must (i) sense all five fingers,
(ii) transmit their data to a phone without a wired tether, (iii) be small and
light enough to wear, and (iv) be buildable from affordable, available parts.
These requirements drove the choice of one IMU per finger, an ESP32-class
microcontroller per finger for local processing and radio, and a phone as the
inference and output device.

\section{Component Selection}

Each finger module pairs a \textbf{Seeed Studio XIAO ESP32-C3} \cite{xiao} with
an \textbf{InvenSense MPU6050} \cite{mpu6050}. The XIAO ESP32-C3 was chosen for
its very small footprint, its integrated Wi-Fi/BLE radio (enabling both ESP-NOW
and BLE), its USB-C interface for flashing and power, and its low cost. The
MPU6050 provides the six axes of inertial data over I\textsuperscript{2}C and is
likewise inexpensive and well documented. The phone runs a cross-platform
Flutter application, so no platform-specific hardware is required on the output
side.

\section{The Finger Module}

A finger module is a compact stack: the XIAO and the MPU6050 are soldered
together with very short wires and separated by an insulating pad that also
limits heat transfer from the microcontroller to the sensor; a second
heat-spreading pad sits on top of the microcontroller and is left exposed to
open air so it can actually dissipate heat rather than trap it. The complete
stack mounts on the back of the fingertip (over the nail). Mounting on the
nail—rather than along the finger segments—keeps the sensor on a rigid part of
the finger, simplifies the wiring, and avoids fitting electronics into the
knuckle gaps.

The I\textsuperscript{2}C bus uses pins D4 (GPIO6, SDA) and D5 (GPIO7, SCL) on
the XIAO at 400\,kHz. Each MPU6050 is configured with a digital low-pass filter
at 10\,Hz, a gyroscope full-scale range of $\pm$500\,$^{\circ}$/s, and an
accelerometer range of $\pm$4\,g.

%% FIGURE: a photo of one nail-mounted XIAO+MPU module.
%% To include: place the photo in Imgs/ (e.g. module.jpg) and uncomment:
% \begin{figure}[!htbp]
%     \centering
%     \includegraphics[width=0.5\textwidth]{module.jpg}
%     \caption{\label{fig:module}A nail-mounted finger module: XIAO ESP32-C3
%     stacked over an MPU6050 with an isolating pad between them.}
% \end{figure}

\section{Network Topology: Master and Slaves}

The five modules form a star network on the hand. The \textbf{middle-finger
module is the master}; the \textbf{thumb, index, ring, and pinky modules are
slaves}. Each slave reads and fuses its own IMU and transmits a small packet to
the master over ESP-NOW at 25\,Hz. The master reads its own IMU, collects the
latest packet from each slave, and assembles a single 45-channel feature vector
(five fingers $\times$ nine fields), which it forwards to the phone over BLE.
This topology is illustrated conceptually in Table~\ref{tab:topology}.

\begin{table}[!htbp]
\centering
\caption{\label{tab:topology}Network roles and data paths.}
\begin{tabular}{lll}
\toprule
Module (finger) & Role & Outgoing link \\
\midrule
Thumb (1)  & Slave  & ESP-NOW $\rightarrow$ master \\
Index (2)  & Slave  & ESP-NOW $\rightarrow$ master \\
Middle (3) & Master & BLE $\rightarrow$ phone \\
Ring (4)   & Slave  & ESP-NOW $\rightarrow$ master \\
Pinky (5)  & Slave  & ESP-NOW $\rightarrow$ master \\
\bottomrule
\end{tabular}
\end{table}

\section{Evolution of the Design}

The final architecture was reached iteratively, and the intermediate designs
are instructive.

\textbf{Iteration 1 — wrist hub with external antennas.} The first design used
a separate ESP32 development board worn at the wrist as a receiver, and clipped
an external U.FL antenna to each finger module to improve the radio link to the
wrist. Measurements revealed three problems: the wrist-to-fingertip distance
made the link sensitive to hand pose (a closed fist shadowed the inner
fingers); the five external antennas, bunched together when the hand closed,
coupled and interfered; and routing the antenna coax along moving fingers was
mechanically fragile.

\textbf{Iteration 2 — nail-mounted modules, no antennas, no hub.} Relocating
the receiver onto the middle fingertip reduced the worst-case link distance
from roughly 20\,cm (wrist) to a few centimetres (adjacent fingertip), which
made the on-chip antennas more than sufficient and removed the external
antennas and their coupling entirely. Because the master fingertip can talk to
the phone directly over BLE, the separate wrist hub was also eliminated. This
single change resolved the antenna-coupling, hand-shadowing, and
mechanical-fragility problems together, and reduced the part count.

\textbf{Sampling rate.} The data rate was reduced from an initial 50\,Hz to
25\,Hz. At 50\,Hz the master could not reliably both receive four ESP-NOW
streams and emit BLE notifications on a single radio; 25\,Hz leaves ample
headroom, and human hand motion is well below the resulting 12.5\,Hz Nyquist
limit, so no useful information is lost.

\section{Power}

For development and evaluation the glove is powered from a regulated 5\,V
supply over USB-C; each XIAO regulates this to 3.3\,V on-board. Power integrity
proved to be a practical constraint: the simultaneous radio transmissions of
several modules produce brief current spikes, and a marginal supply or wiring
path can sag the rail far enough to trigger the ESP32-C3 brown-out detector and
reset a module. The mitigations—reducing the Wi-Fi transmit power to a level
appropriate for a few-centimetre link, improving the power wiring, and adding
local decoupling capacitance—are discussed with the evaluation in Chapter~7.
